\section{Conclusion}


The NTT transforms lattice cryptography from theoretically elegant to practically
deployable. Our Go implementation with AVX2 SIMD achieves within 5\% of the reference C
implementation, demonstrating that Go is a viable language for production cryptographic
code when combined with targeted assembly for inner loops.

The key takeaways:
\begin{enumerate}[nosep]
  \item Precompute twiddle factors in Montgomery form and bit-reversed order.
  \item Use Montgomery multiplication throughout---never divide by $q$.
  \item Vectorize the butterfly with SIMD: 16 butterflies per AVX2 instruction.
  \item Verify with round-trip and homomorphism fuzz tests.
  \item Profile before optimizing: NTT is the bottleneck, not the surrounding code.
\end{enumerate}

\begin{thebibliography}{99}
\bibitem{cooley1965} Cooley, J.W. and Tukey, J.W. ``An Algorithm for the Machine Calculation of Complex Fourier Series.'' Mathematics of Computation, 1965.
\bibitem{montgomery1985} Montgomery, P.L. ``Modular Multiplication Without Trial Division.'' Mathematics of Computation, 1985.
\bibitem{longa2016} Longa, P. and Naehrig, M. ``Speeding up the Number Theoretic Transform for Faster Ideal Lattice-Based Cryptography.'' CANS, 2016.
\bibitem{seiler2018} Seiler, G. ``Faster AVX2 Optimized NTT Multiplication for Ring-LWE Lattice Cryptography.'' IACR ePrint 2018/039.
\bibitem{kyber} Avanzi, R. et al. ``CRYSTALS-Kyber: Algorithm Specifications.'' NIST PQC, 2021.
\bibitem{dilithium} Ducas, L. et al. ``CRYSTALS-Dilithium: Algorithm Specifications.'' NIST PQC, 2021.
\end{thebibliography}

